\documentclass{article}

\usepackage{amsmath,amssymb}
\usepackage{xurl}
\usepackage[hidelinks]{hyperref}
\setlength{\emergencystretch}{2em}

% Shared commands for notes in docs/paper-gaps/.

\DeclareUnicodeCharacter{2080}{\ifmmode{}_{0}\else\textsubscript{0}\fi}
\DeclareUnicodeCharacter{2081}{\ifmmode{}_{1}\else\textsubscript{1}\fi}
\DeclareUnicodeCharacter{2082}{\ifmmode{}_{2}\else\textsubscript{2}\fi}
\DeclareUnicodeCharacter{2099}{\ifmmode{}_{n}\else\textsubscript{n}\fi}

\newcommand{\Lean}{\textsc{Lean}}
\ExplSyntaxOn
\NewDocumentCommand{\leanid}{m}{
  \ifmmode
    \textsf{\detokenize{#1}}
  \else
    \group_begin:
    \urlstyle{sf}
    \tl_set:Nn \l_tmpa_tl {#1}
    \tl_replace_all:Nnn \l_tmpa_tl {\_} {_}
    \exp_args:NV \path \l_tmpa_tl
    \group_end:
  \fi
}
\ExplSyntaxOff
\providecommand{\tr}{\operatorname{tr}}
\newcommand{\tnleanIssueBase}{https://github.com/LionSR/TNLean/issues/}
\newcommand{\tnleanPrBase}{https://github.com/LionSR/TNLean/pull/}
\newcommand{\tnleanDocBase}{https://sirui-lu.com/TNLean/docs/}
\newcommand{\ghissue}[1]{\href{\tnleanIssueBase#1}{GitHub issue \##1}}
\newcommand{\ghpr}[1]{\href{\tnleanPrBase#1}{PR \##1}}
\newcommand{\doclink}[1]{\href{\tnleanDocBase#1.html}{\path{#1}}}

% Dirac notation and the identity operator, as in blueprint/src/macros/common.tex.
% \providecommand keeps note-local definitions authoritative.
\usepackage{dsfont}
\providecommand{\ket}[1]{|#1\rangle}
\providecommand{\bra}[1]{\langle #1|}
\providecommand{\braket}[2]{\langle #1 | #2 \rangle}
\providecommand{\ketbra}[2]{|#1\rangle\!\langle #2|}
\providecommand{\Id}{\mathds{1}}
\providecommand{\one}{\mathds{1}}
\providecommand{\C}{\mathbb{C}}
\providecommand{\MN}[1]{M_{#1}(\C)}
\providecommand{\MD}{M_D(\C)}

% External referencing for paper sources.  The build helper writes sanitized
% auxiliary files under build/paper-aux/ and excludes duplicated source labels.
\usepackage{xr}

% Import a generated paper auxiliary file with a caller-supplied label prefix.
% The candidate paths support builds from docs/paper-gaps/, the repository root,
% and build/paper-aux/ itself.
\newcommand{\importpaperaux}[2]{%
  \IfFileExists{../../build/paper-aux/#2.aux}{%
    \externaldocument[#1]{../../build/paper-aux/#2}%
  }{%
    \IfFileExists{build/paper-aux/#2.aux}{%
      \externaldocument[#1]{build/paper-aux/#2}%
    }{%
      \IfFileExists{#2.aux}{%
        \externaldocument[#1]{#2}%
      }{}%
    }%
  }%
}

% General external reference macros
\newcommand{\paperref}[2]{\ref{#1-#2}}
\newcommand{\papereqref}[2]{(\ref{#1-#2})}

% CPSV16 source (1606.00608 / MPDO-22-12-17-2.tex) external document and shortcuts
\importpaperaux{cpsv16-}{1606.00608/MPDO-22-12-17-2-tnlean}
\newcommand{\cpsvref}[1]{\paperref{cpsv16}{#1}}
\newcommand{\cpsveqref}[1]{\papereqref{cpsv16}{#1}}

% Verdict marker: \gapnote{<kind>}{<status>}. Kind is the policy
% classification (clarification, local-correction, scope-restriction,
% unfaithful, false-source, open-gap); status is open, wip (elimination
% underway), resolved, or historical. Machine-read by the site index;
% prints a verdict line.
\newcommand{\gapnote}[2]{\par\noindent\textbf{Verdict:} #1 (#2).\par}


\title{Schmidt-number premise of Wolf's reduction criterion (Wolf eq.~(3.18))}
\author{}
\date{2026-06-24}

\begin{document}
\maketitle

\section{Source statement}

The reduction-criterion paragraph of Wolf's \emph{Quantum Channels \&
Operations} (\path{Notes/WolfNoteTexSource/ch03_positive_not_completely.tex},
lines 329--338) introduces the maps $T_n(X)=\tr(X)\,\mathbf 1-X/n$, notes that
each is $n$-positive and self-dual, and concludes that
\begin{align}
  n\,\rho_1\otimes\mathbf 1\ge\rho
  \qquad\text{and}\qquad
  n\,\mathbf 1\otimes\rho_2\ge\rho
  \tag{3.18}
\end{align}
is \emph{necessary for a bipartite state $\rho$ to have Schmidt number $n$},
with $\rho_1$ and $\rho_2$ the two reduced density matrices. The source assertion
is the implication
\begin{align}
  \operatorname{sn}(\rho)\le n
  \;\Longrightarrow\;
  \bigl(n\,\rho_1\otimes\mathbf 1\ge\rho\ \text{and}\
  n\,\mathbf 1\otimes\rho_2\ge\rho\bigr),
  \tag{\path{source}}
\end{align}
where $\operatorname{sn}(\rho)$ is the Schmidt number of $\rho$. Its proof has
two steps. A state of Schmidt number at most $n$ is a convex combination of
pure states of Schmidt rank at most $n$, so $n$-positivity of $T_n$ gives
\begin{align}
  \operatorname{sn}(\rho)\le n
  \;\Longrightarrow\;
  (T_n\otimes\operatorname{id})(\rho)\ge 0,
  \tag{step 1}
\end{align}
and the algebraic identity $(T_n\otimes\operatorname{id})(\rho)=\mathbf 1\otimes\rho_2-n^{-1}\rho$
then yields
\begin{align}
  (T_n\otimes\operatorname{id})(\rho)\ge 0
  \;\Longrightarrow\;
  n\,\mathbf 1\otimes\rho_2\ge\rho,
  \tag{step 2}
\end{align}
with the companion inequality obtained from the swapped state.

\section{Formalized statement}

The formal theorems\footnote{Formal declarations:
    \leanid{Matrix.reductionCriterion_left} and
    \leanid{Matrix.reductionCriterion_right} in
    \doclink{TNLean/Channel/ReductionCriterion}.} establish step~2 only. For a
bipartite matrix $\rho$ and $n\ge 1$,
\begin{align}
  (T_n\otimes\operatorname{id})(\rho)\ge 0
  &\;\Longrightarrow\;
  n\,\mathbf 1\otimes\rho_2\ge\rho,\\
  (T_n\otimes\operatorname{id})(\rho^{F})\ge 0
  &\;\Longrightarrow\;
  n\,\rho_1\otimes\mathbf 1\ge\rho,
\end{align}
where $\rho^{F}$ is the image of $\rho$ under the factor swap and the second
hypothesis equals Wolf's symmetric condition $(\operatorname{id}\otimes T_n)(\rho)\ge 0$. The
underlying identity $(T_n\otimes\operatorname{id})(\rho)=\mathbf 1\otimes\rho_2-n^{-1}\rho$ is
also formalized.\footnote{Formal declaration:
    \leanid{Matrix.tensorMapId_tEta_eq}.}

\section{Resolution}

Step~1 is now formalized. A Schmidt-number predicate for mixed
states\footnote{Formal declaration: \leanid{Matrix.HasSchmidtNumberLE} in
    \doclink{TNLean/Channel/SchmidtNumber}.} records that $\rho$ is a finite sum
of pure-state projectors of Schmidt rank at most $n$, which is Wolf's convex
decomposition up to normalization. At $n=1$ this predicate coincides with
separability.\footnote{Formal declaration:
    \leanid{Matrix.hasSchmidtNumberLE_one_iff_isSeparable}.}

The pure-state step shows that for a vector $\psi$ of Schmidt rank at most $n$
(with $1\le n<D$) the ampliation $(T_n\otimes\operatorname{id})(\ketbra\psi\psi)$ is positive
semidefinite,\footnote{Formal declaration:
    \leanid{Matrix.tensorMapId_tEta_posSemidef_of_hasSchmidtRankLE}.} by writing
$\psi$ through the maximally entangled vector as a square right-factor matrix of
rank equal to the Schmidt rank of $\psi$, identifying the ampliation with the
right-factor Choi compression, and using the $n$-positivity of $T_n$ on the
resulting Schmidt-rank-$\le n$ vector. Summing over the pure terms gives step~1:
a state of Schmidt number at most $n$ has
$(T_n\otimes\operatorname{id})(\rho)\ge 0$.\footnote{Formal declaration:
    \leanid{Matrix.HasSchmidtNumberLE.tensorMapId_tEta_posSemidef}.} Composing
step~1 with the formalized step~2 yields Wolf's eq.~(3.18) with its
Schmidt-number premise.\footnote{Formal declarations:
    \leanid{Matrix.reductionCriterion_left_of_hasSchmidtNumberLE} and
    \leanid{Matrix.reductionCriterion_right_of_hasSchmidtNumberLE}.}

The $n$-positivity content used in the pure-state step is the threshold theorem:
the map $T_n$ is $n$-positive for all $1\le n<D$.\footnote{Formal declaration:
    \leanid{Matrix.isNPositiveMap_tEta_iff} in
    \doclink{TNLean/Channel/NPositivityChainStrict}, the reduction map being
    $T_\eta$ at $\eta=n$ by \leanid{Matrix.reductionMap_eq_tEta}.}

\section{Classification}

\emph{Resolved}, with a residual scope restriction $1\le n<D$ inherited from the
$n$-positivity threshold of $T_\eta$, whose top index $n=D$ (complete positivity)
is documented in \path{docs/paper-gaps/wolf_t_eta_top_index_scope.tex}. The full
criterion with Wolf's Schmidt-number premise is now available at
\leanid{Matrix.reductionCriterion_left_of_hasSchmidtNumberLE} and
\leanid{Matrix.reductionCriterion_right_of_hasSchmidtNumberLE}. The operator-half
theorems \leanid{Matrix.reductionCriterion_left} and
\leanid{Matrix.reductionCriterion_right} remain as the separability-free
sub-step, no longer the limit of what is formalized.

\bibliographystyle{alpha}
\bibliography{references}

\end{document}
